% =====================================================================
\chapter{Az eredmények értékelése és döntés: mit építünk be a rendszerbe}
\label{ch:diszkusszio}
\chaptermark{Értékelés és döntés}
% =====================================================================

Ebben a fejezetben \az{\ref{ch:eredmenyek}}. és \az{\ref{ch:modellek}}.~fejezet méréseit értékeljük, és
ezekből döntjük el, mi kerüljön a \Mimir{} éles generálási útvonalába, amely ma a naiv RAG-ot (\arm{E0})
futtatja. Az olvasó megtudja, mely hipotéziseket támasztják alá a mérések, hogyan viszonyulnak
eredményeink az irodalomhoz, és GPU-szintenként melyik konfiguráció a legjobb, a leggyorsabb és a
legrobusztusabb, valamint ebből mi kerül be a rendszerbe (\ref{sec:dontes}.~szakasz). A fejezetet a mérés
tanulságai, az érvényesség korlátai és az etikai szempontok zárják.

% ---------------------------------------------------------------------
\section{A hipotézisek értékelése}
\label{sec:hipotezis-ertekeles}
% ---------------------------------------------------------------------

\begin{table}[htbp]
\centering
\caption[A hipotézisek és a vizsgálatok bizonyítéka.]{A hipotézisek és a vizsgálatok bizonyítéka. Fő: a 30 dokumentumos fő vizsgálat (PULI-LlumiX 8B és
Gemma 3 12B); Előzetes: a hat dokumentumos vizsgálat (Qwen2.5-7B és a szervermodell). $\tilde p$: Holm-korrigált
$p$-érték; G: alátámasztottság; A: vak megoldhatóság.}
\label{tab:hipotezisek}
\small
\begin{tabularx}{\textwidth}{@{}>{\raggedright\arraybackslash}p{2.2cm}>{\raggedright\arraybackslash}X>{\raggedright\arraybackslash}p{2.9cm}@{}}
\toprule
Hipotézis & Bizonyíték & Értékelés \\
\midrule
\ref{h:terv}.\ Tervezés &
Fő: $Q$ $+0{,}26$ (PULI) és $+0{,}23$ (Gemma), $\tilde p<0{,}001$, illetve $0{,}002$; G 27$\to$58, illetve
46$\to$88\,\%; lefedettség 10$\to$23, illetve 24$\to$43\,\%; a forma nem romlott. Előzetes: $+0{,}24$
(6 dok., nem szignifikáns). A rögzített darabolású naiv RAG (\arm{E5c}) a Gemmánál $0{,}97$ &
\textbf{támogatott} a szemantikus darabolású \arm{E0}-hoz képest \\
\ref{h:verif}.\ Ellenőrzés &
Fő: G csökkent (PULI 58$\to$52\,\%, Gemma 88$\to$86\,\%); $Q$ $+0{,}13$ ($\tilde p=0{,}066$), illetve $+0{,}01$
($\tilde p=1{,}00$). Előzetes (Qwen): G 80$\to$90\,\%, A 80$\to$97\,\%, $Q$ $+0{,}08$ (5 dok.\ javult, 0 romlott,
$\tilde p=0{,}19$). $\kappa=0{,}00$--$0{,}08$ minden modellnél &
\textbf{a fő vizsgálatban nem támogatott}; az előzetesben leíró szinten javult \\
\ref{h:lefed}.\ Lefedettség &
Gemma: teljes dok.\ 43\,\% vs.\ \arm{E2} 43\,\%; Qwen: 45 vs.\ 44\,\%; a tömörülés a naiv RAG sajátja (10--25\,\%) &
\textbf{nem támogatott} \\
\ref{h:helyi}.\ Helyi vs.\ szerver &
Előzetes, közös dokumentumokon: Qwen \arm{E2} $0{,}95$ vs.\ szerver \arm{E4} $0{,}93$ (három angol dok.), \arm{E2}
$0{,}96$ vs.\ \arm{B-doc-S} $0{,}95$ (hat dok.). A szervermodell a saját módszerünkkel (\arm{E4b}) $0{,}99$. A hat
dokumentumon a Gemma \arm{E2} $0{,}97$, a PULI \arm{E2} $0{,}72$. A modellek 4,8--7,8\,GiB-ot foglalnak &
\textbf{támogatott} (leíró, modellfüggő) \\
\ref{h:meres}.\ Mérés &
Az oktatói értékelés folyamatban (\ref{sec:oktatoi-eredmeny}.~szakasz) &
\textbf{nyitott} \\
\bottomrule
\end{tabularx}
\end{table}

Megjegyzések a táblázathoz. \textbf{\ref{h:terv}.} Az \arm{E1} és
az \arm{E0} több tényezőben különbözik (lekérdezés, kérdésenkénti generálás), ezért a javulást a vizsgaterv
egészének tulajdonítjuk, nem egyetlen lépésének. \textbf{\ref{h:verif}.} A lánc a vizsga egészét javította
(forma, szivárgás, nem kérdő törzsek; \ref{sec:fo-ellenorzo}.~szakasz), a kérdések tartalmi helyességét nem.
\textbf{\ref{h:lefed}.} A PULI teljes dokumentumos vizsgáinak alacsony lefedettsége a kevés kérdésből
adódik, és tíz kérdésnél a hosszú dokumentumokon a felső korlát minden konfigurációt egyformán korlátoz.
\textbf{\ref{h:helyi}.} A hipotézis a szervermodell \emph{naiv} változatához mér; a „helyi modell elérte a
szervermodellt” állítás csak a szerveres teljes dokumentumos promptra igaz, a saját módszerünkkel futó
szervermodellre nem. \textbf{\ref{h:meres}.} Addig a bíró ítéleteit relatív, a konfigurációk összevetésére
szolgáló mérésként kezeljük. \tbd{[Az oktatói értékelés eredménye után: a hipotézis értékelése a
Krippendorff-$\alpha$ alapján.]}

% ---------------------------------------------------------------------
\section{Összevetés az irodalommal}
\label{sec:irod-osszevetes}
% ---------------------------------------------------------------------

Az alábbiakban eredményeinket \az{\ref{ch:szakirodalom}}.~fejezetben tárgyalt munkákkal vetjük össze. Az
irodalmi eredményeket minőségileg idézzük, mert a feladatok, a modellek és a mutatók eltérnek a mieinktől;
az összevetés ezért irányt és nem pontos nagyságot jelez.

\paragraph{Önellenőrzés külső visszajelzés nélkül.} Huang és munkatársai~\cite{huang2024selfcorrect} szerint
a modellek külső visszajelzés nélkül nem tudják megbízhatóan javítani saját érvelésüket, és a
teljesítmény az önjavítás után akár romolhat is. Eredményeink ezzel összhangban vannak: a PULI-nál az
ellenőrző lánc haszna lényegében a szabályokból jött (a csak szabályalapú \arm{E2x} ugyanannyit ért el, mint
az \arm{E2}: $0{,}67$ és $0{,}67$), a Gemmánál a különbség nem szignifikáns, a modellalapú ítélet pedig minden
modellnél a véletlen szintjén egyezett a bíróval. A Qwennél az \arm{E2x} rosszabb volt ($0{,}89$ vs.\
$0{,}96$, egy dokumentumon jobb, ötön rosszabb), ami arra utal, hogy a forrásszöveggel szembeni ítélet --
vagyis külső információ -- egy közepes modellnél adhat hozzá. Ezt az \arm{E2x} hibás futása miatt újra kell
mérni (\ref{sec:dontes}.~szakasz).

\paragraph{Darabolás.} Qu, Tu és Bao~\cite{qu2024chunking} szerint a szemantikus darabolás többletköltségét
a mért előnyök nem indokolják. Nálunk a rögzített, 800 karakteres darabolás (\arm{E5c}) a Gemmánál
($0{,}72\to0{,}97$, 23 dokumentumon jobb) és a Qwennél ($0{,}64\to0{,}86$, leíró szinten) is jobb volt, mint a
szemantikus \arm{E0}; a PULI-nál egyik változat sem adott értékelhető vizsgát. Eredményünk a
megállapításukat egy konkrét mechanizmussal egészíti ki: a minimális darabméret nélküli szemantikus
darabolás rövid, tartalom nélküli darabokat hoz létre, és a feladatleírás típusú lekérdezés éppen ezeket
találja meg (\ref{sec:fo-robusztus}.~szakasz).

\paragraph{Lekérdezés-átírás.} A HyDE~\cite{gao2023hyde} abból indul ki, hogy a lekérdezésnek a keresett
szöveghez kell hasonlítania, ezért egy generált hipotetikus dokumentum beágyazásával keres. A mi
eredményünk ugyanezt támasztja alá a vizsgagenerálásban: a kereső tartalmi kérdéssel jól működik (a
referenciakérdések igazoló darabja a fő vizsgálat dokumentumain 86\,\%-ban az első három találat között volt), a naiv
RAG kudarca tehát a feladatleírásból képzett lekérdezésből fakad. A vizsgaterv tartalmi fogalmai a
lekérdezés-átírás egy egyszerű formájának tekinthetők; egy HyDE-szerű lekérdezés a naiv útvonalon
továbbfejlesztési lehetőség.

\paragraph{Hosszú kontextus és RAG.} Li és munkatársai~\cite{li2024ragvslc} szerint elegendő erőforrás
mellett a hosszú kontextus átlagosan jobb a RAG-nál, a RAG viszont jóval olcsóbb, a lekérdezések 63\,\%-án a
kettő előrejelzése azonos, és a kérdésenként választó útválasztással a költség csökkenthető. Nálunk ez modellfüggő: a Gemma teljes
dokumentumos promptja ($0{,}96$) versenyképes a tervezett folyamattal és a jól darabolt naiv RAG-gal, a
gyengébb modelleknél viszont a teljes dokumentumos prompt elmarad (Qwen $0{,}83$, PULI $0{,}41$). Az
útválasztás gondolata a mi adaptív ellenőrzésünkkel rokon (\ref{sec:dontes}.~szakasz).

\paragraph{A hosszú bemenet egyenetlen használata.} Liu és munkatársai~\cite{liu2024lost} eredménye alapján
azt vártuk, hogy a teljes dokumentumot kapó modell kérdései a dokumentum egyes részeire tömörülnek
(\ref{h:lefed}.~hipotézis). A lefedettségben ez nem jelent meg. Leíró szinten a Gemma teljes dokumentumos
promptja a hosszú dokumentumokon valamivel gyengébb ($0{,}94$), mint a rövideken ($0{,}98$), ami a hossz
növelésével romló teljesítmény irányába mutat~\cite{hsieh2024ruler}, de a különbség kicsi, és a bíró 24\,000
karakteres ablaka miatt a hosszú dokumentumokon az alátámasztottság alsó becslés.

\paragraph{Hibás kulcsok a generált kérdésekben.} Doughty és munkatársai~\cite{doughty2024comparative}
szerint a GPT-4 kérdéseinek 81,7\,\%-a minden szempontnak megfelelt, de 4,9\,\%-ukban több helyes válasz
volt (az oktatói kérdéseknél 1,1\,\%), 4,0\,\%-ukban pedig egy disztraktor elárulta a kulcsot (0,9\,\%); Wang és munkatársai~\cite{wang2022towards} a GPT-3
kérdéseinek egy részét az emberiekkel összemérhetőnek találták. A mi méréseink ugyanezt a képet adják: a
Gemma tervezett konfigurációiban a bíró a kérdések 85--88\,\%-át használná fel, a Qwen \arm{E2} kérdéseinek
80\,\%-át, a szervermodell \arm{E4b} kérdéseinek 98\,\%-át, és még a legjobb helyi konfigurációk kérdéseinek
is 4--14\,\%-a a bíró szerint nem alátámasztott. A mutatók és a tárgyterületek eltérnek, így ez nem közvetlen
összevetés, de a következtetés azonos: oktatói átnézés nélkül hibás kérdések kerülnének a vizsgába.

\paragraph{A bíró érvényessége.} Zheng és munkatársai~\cite{zheng2023judge} szerint egy erős bíró több mint
80\,\%-ban egyezik az emberi preferenciákkal, nagyjából annyira, mint az emberek egymással, és pozíciós,
hossz- és önpreferencia-torzítást is azonosítottak; Chiang és Lee~\cite{chiang2023human} az emberi
szakértőkével azonos utasítással kapott hasonló egyezést. Ezek az eredmények páros preferenciákra és más
feladatokra vonatkoznak, ezért nem vihetők át a mi forrásalapú, pontszámos rubrikánkra. A vizsgálatok több
bírói vakfoltot is feltártak: a szivárgó és a tartalmilag üres kérdéseket a bíró magasra pontozta, a teljes
dokumentumos kérdéseket a 24\,000 karakteres ablak miatt alulpontozhatta, és ítéletei forrásfüggően
hiányoztak (\ref{sec:meres-tanulsag}.~szakasz). A más modellcsaládba tartozó bíró az önpreferenciát
csökkenti~\cite{panickssery2024selfpref}, de a bíró érvényességét csak az oktatói validálás döntheti el
(\ref{h:meres}.~hipotézis).

% ---------------------------------------------------------------------
\section{Mikor éri meg a módszer?}
\label{sec:mikor}
% ---------------------------------------------------------------------

A módszer haszna a modell képességétől függ (\ref{fig:lepesek}.~ábra). Gyenge formátumkövető modellnél
(PULI-LlumiX 8B) csak a felbontás ad értékelhető vizsgát, de a modellt ez sem teszi alkalmassá. Közepes
modellnél (Qwen2.5-7B) éri meg a legjobban: a többlethívások árán a szerveres teljes dokumentumos prompt
szintjét éri el, összhangban azzal, hogy a futás közbeni többletszámítás a nagyobb modellt részben
helyettesítheti~\cite{snell2025scaling}. Erős modellnél (Gemma 3 12B) átlagos minőségben nem ad hozzá, értéke
a kérdésenkénti forrás (P1) és a terv (P2). A szintenkénti számokat \az{\ref{tab:dontes}}.~táblázat adja.

% ---------------------------------------------------------------------
\section{Döntés: mit építünk be}
\label{sec:dontes}
% ---------------------------------------------------------------------

\Az{\ref{tab:dontes}}.~táblázat GPU-szintenként a legjobb minőségű, a leggyorsabb és a legrobusztusabb
(a leggyengébb dokumentumon is legjobb) konfigurációt mutatja.

\begin{table}[htbp]
\centering
\caption[Döntési táblázat GPU-szintenként.]{Döntési táblázat GPU-szintenként. $Q$: minőségi index; Min.: a leggyengébb dokumentum
$Q$-ja; Idő: a medián futási idő szorzója az azonos modellel futó \arm{E0}-hoz képest. A 8\,GB-os szint
a hat dokumentumos előzetes, a 12\,GB-os a 30 dokumentumos fő vizsgálat adatai. A szintén 8\,GB-on futó
PULI-LlumiX 8B egyik szempontból sem ajánlott (legjobb \arm{E2} $0{,}67$, a leggyengébb dokumentumon
$0{,}38$, 18 dokumentum $0{,}70$ alatt).}
\label{tab:dontes}
\small
\setlength{\tabcolsep}{4pt}
\begin{tabularx}{\textwidth}{@{}>{\raggedright\arraybackslash}p{2.6cm}lccc>{\raggedright\arraybackslash}X@{}}
\toprule
Szempont & Konfig. & $Q$ & Min. & Idő & Megjegyzés \\
\midrule
\multicolumn{6}{@{}l}{\textit{8\,GB: Qwen2.5-7B}} \\
Legjobb minőség & \arm{E2} & 0,96 & 0,89 & $\approx11\times$ & a szerveres teljes dokumentumos prompt
($0{,}95$) szintje; kérdésenkénti forrás \\
Leggyorsabb & \arm{E5c} & 0,86 & 0,50 & $\approx1\times$ & a saját magyar szövegeken $0{,}77$, a törzsek
44\,\%-a nem kérdés; csak vázlatnak \\
Legrobusztusabb & \arm{E2} & 0,96 & 0,89 & $\approx11\times$ & azonos a legjobb minőségűvel; a következő
az \arm{E2h} (0,88) \\
\textbf{Beépül} & \arm{E2} & 0,96 & 0,89 & $\approx11\times$ & adaptív ellenőrzéssel (lent) \\
\midrule
\multicolumn{6}{@{}l}{\textit{12\,GB: Gemma 3 12B}} \\
Legjobb minőség & \arm{E5c} & 0,97 & 0,81 & $\approx1\times$ & a \arm{B-doc-L}, \arm{E2f}, \arm{E2h}, \arm{E5a}
($0{,}96$) a konfidencia-intervallumon belül; nincs kérdésenkénti forrás és terv \\
Leggyorsabb & \arm{E5c} & 0,97 & 0,81 & $\approx1\times$ & az \arm{E0}, \arm{E5a} és \arm{B-doc-L} is
$\approx1\times$ \\
Legrobusztusabb & \arm{E2h} & 0,96 & 0,88 & $\approx8\times$ & a tervezett konfigurációk minimuma
$0{,}80$--$0{,}88$, egy-egy dokumentumon múlik; a hibrid keresés más modelleknél rontott \\
\textbf{Beépül} & \arm{E1}--\arm{E2} & 0,95 & 0,80--0,82 & $\approx5$--$8{,}5\times$ & tervezett folyamat adaptív
ellenőrzéssel; gyors vázlathoz \arm{E5c} \\
\bottomrule
\end{tabularx}
\end{table}

A 8\,GB-os szinten a tervezett és ellenőrzött folyamat egyszerre a legjobb és a legrobusztusabb. A
12\,GB-os szinten a győztesek minősége a konfidencia-intervallumokon belül azonos, ezért azt választjuk, amit
csak a tervezett folyamat ad meg: a kérdésenkénti forrást és a tervet; a rögzített darabolású naiv RAG gyors
vázlatként marad meg.

\subsection*{Mi kerül be}

\begin{description}[style=nextline,leftmargin=0.6cm]
\item[Kérdésenkénti forráshivatkozás (P1).] Még a legjobb helyi konfigurációk kérdéseinek egy része is
alátámasztatlan (\ref{sec:irod-osszevetes}.~szakasz), és a bíró sem validált, ezért az oktatói ellenőrzés nem
hagyható el. A forrásdarab megjelenítése ezt a szerkesztőben egy kattintásra rövidíti, és a hivatkozás az
automatikus alátámasztottsági ellenőrzés feltétele is.
\item[Vizsgaterv (P2).] A terv kikényszeríti a darabszámot, a típust és a Bloom-kiosztást, ami egy
promptban csak kérés: a PULI egy hívásban átlagosan 1,8 kérdést adott a kért tíz helyett, a tervvel 9,9-et. A Gemmánál a bíró szerint megvalósult magasabb rendű kérdések aránya az
egyhívásos konfigurációk 1--4\,\%-áról 7--11\,\%-ra nőtt; ez is jóval elmarad a tervezett 43\,\%-tól, amit a
felületen jelezni kell (\ref{sec:fo-jelleg}.~szakasz).
\item[Szabályalapú ellenőrzés.] A forma, a szivárgás, a kérdésforma és a duplikátum szabályai olcsók,
determinisztikusak, és a lánc hasznának nagy részét adják (\arm{E2x} $\approx$ \arm{E2}, lásd az
önellenőrzésről szóló bekezdést). A szivárgás szabályát a pilot tette szükségessé
(\ref{sec:szivargas}.~szakasz).
\item[Rögzített vagy minimális méretű darabolás a naiv útvonalon is.] A legolcsóbb mért javítás: a Gemmánál a
$0{,}70$ alatti dokumentumok száma 16-ról nullára csökkent, változatlan futási idő mellett. A minimális
darabméretű szemantikus darabolást és a tervezett folyamat darabolásfüggését nem mértük, ezek elvárások.
\item[Adaptív ellenőrzés.] Az elsőre elfogadott kérdéshelyek aránya a Gemmánál 70\,\%, a Qwennél 23\,\%, a
PULI-nál 5\,\%; minél kisebb ez az arány, annál nagyobb az ellenőrzés haszna
(\ref{tab:hipotezisek}.~táblázat), költsége viszont minden modellnél jelentős. A rendszer ezért az első kérdéshelyeken a teljes láncot futtatja, és ha
az első kísérletek sikeraránya magas, a továbbiakon csak a szabályokat. A küszöb még meghatározandó
(\tbd{[küszöb: a bevezetés előtt mérendő]}); a változatot nem mértük, hatása elvárás.
\end{description}

\subsection*{Mi marad ki}

\begin{description}[style=nextline,leftmargin=0.6cm]
\item[Fogalomgráf (\arm{E3}).] Egyik modellnél sem javított (Gemma $-0{,}01$, PULI $-0{,}10$, Qwen $-0{,}07$),
miközben a legdrágább tervezett változat (Gemma: $\approx10{,}6\times$, 55,5 hívás az \arm{E2} 43 hívásával szemben).
\item[Hibrid keresés (\arm{E2h}).] A Gemmánál nem szignifikánsan jobb ($+0{,}01$), a PULI-nál ($-0{,}08$, 17
dokumentumon rosszabb) és a Qwennél ($-0{,}04$, öt dokumentumon rosszabb) rontott, és külön lexikális indexet
igényel; a Gemmánál mért legmagasabb minimuma egyetlen dokumentumon múlik.
\item[Az LLM-alapú ellenőrzés mint kérdésenkénti ítélet.] Ítélete a bíróéval a véletlen szintjén egyezik; a
Qwennél a bíró szerint alátámasztott kérdések közel felét is megjelölte, a Gemmánál szinte mindent
átengedett. Figyelmeztetésként („ezt nézd meg alaposabban”) megmarad, de a felület nem jelöli a kérdést
„ellenőrzöttnek”, mert ez az automatizálási torzítás~\cite{parasuraman2010complacency} miatt éppen az
átengedett kérdések átnézését gyengítené.
\end{description}

\paragraph{Modellválasztás.} Strukturált feladatban az utasításkövetés fontosabb a magyar nyelvi
specializációnál -- ezt a PULI-LlumiX 8B esetére szűkítve állítjuk, amely a modellkártya szerint a
Llama-3.1-8B-Instruct utólagos utasítás-hangolás nélküli magyar továbbelőtanítása, és amelynél a többnyelvű
Gemma 3 12B minden konfigurációban, a Qwen2.5-7B a hat közös dokumentumon jobb volt
(\ref{ch:modellek}.~fejezet). A rendszer 8\,GB-on a Qwen2.5-7B-t, 12\,GB-on a Gemma 3 12B-t ajánlja.

\paragraph{Fejlesztési ütemterv} (időpontok: \tbd{[ütemezés a témavezetővel egyeztetve]}).
(1)~Rögzített darabolás a meglévő naiv útvonalon. (2)~A tervezett folyamat (vizsgaterv, forráshivatkozás,
szabályalapú és adaptív ellenőrzés) integrálása a Bifrost háttérfeladatába, GPU-szint szerinti
modellválasztással. (3)~A v1 szerkesztő éles backendje az ál-backend helyett. (4)~Ezekkel párhuzamosan a
hiányzó mérések: a hiányzó bírói ítéletek újrapontozása, az \arm{E2x} újramérése működő második
ellenőrzővel, az oktatói validálás (\ref{h:meres}.~hipotézis), a szervermodell a 30 dokumentumon, valamint az
adaptív küszöb és a minimális darabméret mérése.

% ---------------------------------------------------------------------
\section{Mit tanultunk az ellenőrzésről és a mérésről?}
\label{sec:meres-tanulsag}
% ---------------------------------------------------------------------

Az ellenőrző lánc szűrőként és „második esélyként” működik, nem osztályozóként: jelzései egyenként
zajosak, de az elutasított kérdés helyére írt új kérdés átlagosan jobb. A mérés több ponton csendben
torzított volna (tartalékvizsga, csonkolt kódoló, szivárgás, a bíró ablaka, a hibás \arm{E2x}-hívások), miközben
hihető számokat adott; a nyers kimenetek kézi átolvasása ezért nem hagyható el. A bírói ítéletek
forrásfüggően hiányoznak: a Wikipédia-cikkek kérdései kisebb arányban kaptak ítéletet (a Gemma \arm{E2}-ben
összesen 52\,\%, a PULI \arm{E2}-ben 29\,\%). Ez a következtetéseket nem fordítja meg (\ref{sec:fo-meres}.~szakasz),
de a forrásonkénti összevetéseket gyengíti.

% ---------------------------------------------------------------------
\section{Az érvényességet veszélyeztető tényezők}
\label{sec:validitas}
% ---------------------------------------------------------------------

\paragraph{Belső.} Az \arm{E1} és az \arm{E0} több tényezőben különbözik, az ellenőrzés a kérdések hosszát is
megváltoztatja, és a tervezett konfigurációk egy maggal futottak, így a dokumentumonkénti értékekben futási
zaj is van. A vizsgálatok közötti összevetés minden esetben azonos dokumentumokon, a hat közös
dokumentumon történt. Az \arm{E2x} a tervezettől eltérő konfigurációt mért.
\paragraph{Konstrukciós.} A formai megfelelés vizsgaszintű, ezért minden komponenst külön is közlünk. A
tartalmi mutatók egyetlen, még nem validált bíróra épülnek, amely csak a kérdés saját forrását látja, és
ítéletei forrásfüggően hiányoznak (\ref{sec:fo-meres}.~szakasz).
\paragraph{Külső.} A 30 dokumentumból 24 magyar Wikipédia-cikk, az angol részhalmaz három dokumentum; a
szervermodell csak a hat dokumentumon futott; csak feleletválasztós kérdéseket értékeltünk; a nyílt
dokumentumok a tanítóadatban is szerepelhettek~\cite{sainz2023contamination}. A futási időt ugyanazon a GPU-n
mért relatív szorzóként értelmezzük.
\paragraph{Következtetési.} A 30 dokumentum a nagy hatások kimutatására elegendő, a kicsikére (például az
ellenőrzés hatására a Gemmánál) nem; a hat dokumentumos vizsgálat következtetései leíró jellegűek.

% ---------------------------------------------------------------------
\section{Etikai és jogi szempontok}
% ---------------------------------------------------------------------

A döntés a jogi keret (\ref{sec:irod-jog}.~szakasz) követelményeit is szolgálja: a beépülő változat helyben
futtatható, így helyi módban (\code{LOCAL\_ONLY}) a tananyag nem hagyja el az oktató gépét; a feltöltött
tartalom a feladat végén törlődik (\ref{sec:zero}.~szakasz); a forráshivatkozás és a kötelező oktatói
jóváhagyás az emberi felügyeletet konkrét munkafolyamattá teszi. A rendszer kérdést generál, nem értékeli a
hallgatókat, és az export feltünteti az MI-közreműködést~\cite{aiact2024}.

\paragraph{Tanulság.} A \Mimir{} a naiv RAG helyett a tervezett folyamatot építi be forráshivatkozással,
szabályalapú és adaptív ellenőrzéssel, a naiv útvonalon rögzített darabolással; a fogalomgráf és a hibrid
keresés kimarad. Az erős modellnél a döntés az ellenőrizhetőségen alapul, ezért az oktatói validálás
erősítheti meg vagy módosíthatja.
